% This is a simple sample document.  For more complicated documents take a look in the exercise tab. Note that everything that comes after a % symbol is treated as comment and ignored when the code is compiled.

\documentclass{article} % \documentclass{} is the first command in any LaTeX code.  It is used to define what kind of document you are creating such as an article or a book, and begins the document preamble

\usepackage{amsmath} % \usepackage is a command that allows you to add functionality to your LaTeX code

\title{Simple Sample} % Sets article title
\author{My Name} % Sets authors name
\date{\today} % Sets date for date compiled

% The preamble ends with the command \begin{document}
\begin{document} % All begin commands must be paired with an end command somewhere
  \maketitle % creates title using information in preamble (title, author, date)
  


  % Solve these questions in "hw02.tex" based on James K. Strayer's Elementary Number Theory contexts with the rules as follows.
  % 1. Use existing questions and answers to have consistency in proof style and phrasing (passive voice, 3rd-person excluding "the one", "we" ,...)
  % 2. Bring references from the selected lines of "_.tex" when first used, then cross-reference with \refname{}, but not too repetitive. Like existing questions and answers.
  % 3. Remind "Goldilocks principle" along with the existing content: be concise but logically complete. Omit obvious steps, but show all crucial reasoning needed for verification. For example: "To show $f(n)=n+1$ is bijective: if $f(a)=f(b)$ then $a+1=b+1$ so $a=b$, and for any $y$ take $x=y-1$ so $f(x)=y$."

  \newpage


 
  Goldilocks principle - super short \& simple
  One-line summary: write just enough to let the reader check your reasoning - no missing steps, no obvious repetition.

  Quick 3-step checklist to apply it:
  1. State clearly what you will prove.
  2. Give only the definitions/steps needed to connect $\mathrm{A} \rightarrow \mathrm{B}$.
  3. Remove anything the reader can easily fill in, add anything they can't.

  Short example sentence (just-right):
  "To show $f(n)=n+1$ is bijective: if $f(a)=f(b)$ then $a+1=b+1$ so $a=b$, and for any $y$ take $x=y-1$ so $f(x)=y . "$

\end{document} % This is the end of the document


